% Préambule commun aux manuels de Martinet (français et anglais).
% Compilation : pdflatex manuel-fr.tex (deux fois), idem manuel-en.tex, depuis docs/src/.
\documentclass[11pt,a4paper]{article}
\usepackage[utf8]{inputenc}
\usepackage[T1]{fontenc}
\usepackage{lmodern}
\usepackage{microtype}
\usepackage[a4paper,margin=2cm,headheight=24pt]{geometry}
\usepackage{graphicx}
\usepackage[table]{xcolor}
\usepackage[most]{tcolorbox}
\usepackage{enumitem}
\usepackage{booktabs}
\usepackage{array}
\usepackage{tabularx}
\usepackage[nobottomtitles*]{titlesec}
\usepackage{needspace}
\usepackage{fancyhdr}
\usepackage{fontawesome5}
\usepackage{etoolbox}
\DeclareUnicodeCharacter{2192}{\ensuremath{\rightarrow}}
\DeclareUnicodeCharacter{2191}{\ensuremath{\uparrow}}
\DeclareUnicodeCharacter{2193}{\ensuremath{\downarrow}}
\DeclareUnicodeCharacter{25B2}{\ensuremath{\blacktriangle}}
\DeclareUnicodeCharacter{25BC}{\ensuremath{\blacktriangledown}}
\DeclareUnicodeCharacter{21BA}{\ensuremath{\circlearrowleft}}
\DeclareUnicodeCharacter{21BB}{\ensuremath{\circlearrowright}}
\usepackage{amssymb}

% Palette de l'application : bleu ciel, rouge brique pour les alertes
\definecolor{bordeaux}{HTML}{14548C}
\definecolor{bordeauxl}{HTML}{E6F0F8}
\definecolor{alert}{HTML}{B5382B}
\definecolor{alertl}{HTML}{F8E6E3}
\definecolor{or}{HTML}{B8860B}
\definecolor{orl}{HTML}{FBF3DC}
\definecolor{sapin}{HTML}{2C7A55}
\definecolor{sapinl}{HTML}{E2F1E9}
\definecolor{sky}{HTML}{1C6DB4}
\definecolor{skyl}{HTML}{E6F0F8}
\definecolor{summit}{HTML}{6B7A8F}
\definecolor{night}{HTML}{2A2420}
\definecolor{ink}{HTML}{1F2430}
\AtBeginDocument{\color{ink}}

\titleformat{\section}
  {\Large\bfseries\sffamily\color{bordeaux}}
  {\makebox[1.6em][l]{\thesection}}{0.3em}{}
  [{\vspace{1pt}\color{bordeaux!35}\titlerule[1.3pt]}]
\titlespacing*{\section}{0pt}{14pt}{7pt}
\titleformat{\subsection}{\large\bfseries\sffamily\color{night}}{\thesubsection}{0.5em}{}
\pretocmd{\section}{\needspace{10\baselineskip}}{}{}
\pretocmd{\subsection}{\needspace{7\baselineskip}}{}{}
\widowpenalty=10000
\clubpenalty=10000
\raggedbottom
\setlength{\parindent}{0pt}
\setlength{\parskip}{6pt}
\setlist[enumerate]{itemsep=3pt,leftmargin=6mm}
\setlist[itemize]{itemsep=2pt,leftmargin=5mm}

\pagestyle{fancy}
\fancyhf{}
\renewcommand{\headrulewidth}{0.4pt}
\renewcommand{\headrule}{\hbox to\headwidth{\color{bordeaux!40}\leaders\hrule height \headrulewidth\hfill}}
\fancyfoot[C]{\sffamily\small\color{summit}\thepage}
\renewcommand{\sectionmark}[1]{\markboth{#1}{}}

\newcommand{\menu}[1]{\textsf{\textbf{\textcolor{bordeaux}{#1}}}}
\newcommand{\touche}[1]{\tcbox[on line,boxsep=0pt,left=2pt,right=2pt,top=1.5pt,bottom=1.5pt,
  colback=white,colframe=summit,boxrule=0.5pt,arc=1.5pt]{\footnotesize\textsf{#1}}}
\newcommand{\fichier}[1]{\texttt{\small #1}}

\newcommand{\boite}[5]{\newtcolorbox{#1}{enhanced,unbreakable,colback=#2,colframe=#3,
  coltitle=white,fonttitle=\bfseries\sffamily\small,boxrule=0.9pt,arc=2.5pt,
  left=2.5mm,right=2.5mm,top=1.2mm,bottom=1.2mm,title={#4~~#5}}}
\newcommand{\capture}[2][1]{\par\needspace{8\baselineskip}\begin{center}\setlength{\fboxsep}{0pt}\setlength{\fboxrule}{0.6pt}%
  \fcolorbox{summit!50}{white}{\includegraphics[width=\dimexpr#1\linewidth-1.3pt\relax]{#2}}\end{center}}

\usepackage[french]{babel}
\usepackage[colorlinks=true,linkcolor=bordeaux,urlcolor=bordeaux]{hyperref}
\hypersetup{pdftitle={Martinet — manuel d'utilisation},pdfauthor={},pdfcreator={Martinet}}
\boite{astucebox}{orl}{or!85!black}{\faLightbulb}{Bon à savoir}
\boite{securitebox}{sapinl}{sapin}{\faShield*}{Ce que Martinet ne fait jamais}
\boite{attentionbox}{alertl}{alert}{\faExclamationTriangle}{Attention}
\newcommand{\astuce}[1]{\par\needspace{5\baselineskip}\begin{astucebox}#1\end{astucebox}}
\newcommand{\attention}[1]{\par\needspace{5\baselineskip}\begin{attentionbox}#1\end{attentionbox}}
\fancyhead[L]{\sffamily\small\color{bordeaux}\textbf{MARTINET — MANUEL}}
\fancyhead[R]{\sffamily\small\color{summit}\leftmark}
\input{version}

\begin{document}

\begin{titlepage}
\centering
\vspace*{2.2cm}
\includegraphics[width=3.6cm]{../../logo_512.png}\par
\vspace{8mm}
{\sffamily\Huge\bfseries\color{bordeaux} MARTINET\par}
\vspace{3mm}
{\sffamily\LARGE\color{ink} Tenir Windows à jour d'un seul geste\par}
\vspace{2mm}
{\sffamily\large\color{summit} Manuel d'utilisation \textbullet{} version \versionmartinet\par}
\vspace{12mm}
\begin{tcolorbox}[enhanced,width=0.82\linewidth,colback=skyl,colframe=sky,arc=4pt,boxrule=1.2pt,left=5mm,right=5mm,top=3mm,bottom=3mm]
\sffamily\centering
Le martinet passe sa vie en vol et ne se pose presque jamais :\\
Martinet fait sa tournée tout seul.\\[2mm]
\raggedright
\textbf{Tout vérifier} \quad → \quad \textbf{Mettre à jour et nettoyer} \quad → \quad \textbf{Planification}
\end{tcolorbox}
\vfill
{\footnotesize\color{summit}\sffamily Windows 10 et 11\par}
\end{titlepage}

\tableofcontents
\newpage

\section{À quoi sert Martinet}

Sous Linux, une seule commande (\texttt{apt}, \texttt{yay}…) met tout à jour. Sous Windows, chaque
programme a sa propre mise à jour, les pilotes en ont une autre, Windows encore une autre. Martinet
réunit tout cela dans une fenêtre :

\begin{center}
\renewcommand{\arraystretch}{1.35}
\begin{tabularx}{\linewidth}{>{\raggedright\arraybackslash}X >{\sffamily\bfseries\color{bordeaux}}l c}
\toprule
\textbf{Je veux…} & \textbf{\color{ink}Page} & \textbf{Section} \\
\midrule
savoir en un coup d'œil ce qui est en retard & Tableau de bord & \ref{sec:accueil} \\
mettre à jour mes logiciels sans les ouvrir un par un & Programmes (winget) & \ref{sec:programmes} \\
suivre aussi les logiciels que winget ne connaît pas & Hors winget & \ref{sec:hors} \\
installer Windows Update, les pilotes, le pilote de la carte graphique & Windows et pilotes & \ref{sec:windows} \\
retrouver les raccourcis qui ne mènent plus nulle part & Raccourcis morts & \ref{sec:raccourcis} \\
récupérer de la place sur le disque & Ménage & \ref{sec:menage} \\
que tout cela se fasse sans moi & Planification & \ref{sec:planif} \\
\bottomrule
\end{tabularx}
\end{center}

\begin{securitebox}
\begin{itemize}
\item Il ne « nettoie » pas le registre : aucun gain mesurable, un vrai risque de casser un logiciel.
\item Il n'efface aucun raccourci : ils sont \textbf{mis de côté} et se restaurent d'un clic.
\item Il ne touche à aucun fichier personnel ; les caches des navigateurs ne sont vidés que si vous les cochez.
\item Il ne télécharge un installeur que depuis l'adresse \textbf{officielle} de l'éditeur, en HTTPS.
\item Il n'interroge jamais un lecteur réseau : un NAS en veille reste en veille.
\end{itemize}
\end{securitebox}

\section{Installer et lancer}

Téléchargez \fichier{Martinet-Setup.exe} sur la page \textbf{Releases} du projet et lancez-le.
L'installation se fait dans votre compte, sans droits administrateur. Martinet se met ensuite
lui-même à jour : il le propose au démarrage quand une nouvelle version est publiée.

Au lancement, Windows demande \textbf{une seule fois} l'autorisation administrateur : Windows Update,
une partie du ménage et la planification en ont besoin. Si vous refusez, Martinet fonctionne quand
même ; un bandeau jaune rappelle ce qui manque et propose de relancer avec les droits.

Au tout premier lancement, Martinet propose d'activer le \textbf{passage automatique} chaque dimanche
à 12 h (section~\ref{sec:planif}).

\section{Le tableau de bord}\label{sec:accueil}

\capture{../capture_accueil-fr.png}

\begin{enumerate}
\item \menu{Tout vérifier} interroge winget, les éditeurs, Windows Update, NVIDIA et AMD, cherche les
raccourcis morts et mesure le ménage. \textbf{Rien n'est modifié} pendant la vérification, qui prend
une à trois minutes (Windows Update est le plus lent).
\item Chaque carte résume une page ; un clic sur la carte l'ouvre.
\item \menu{Mettre à jour et nettoyer} fait tout ce qui peut se faire sans vous : programmes winget,
mises à jour Windows (facultatives comprises si vous l'avez choisi), signatures de Defender, ménage
coché, et, selon les choix de la page Planification, pilote NVIDIA et raccourcis morts mis de côté.
Avant de commencer, Martinet montre la liste exacte, avec les nombres, et ce qu'il ne fera pas.
\end{enumerate}

\astuce{Le bouton \menu{Afficher le détail}, en bas à droite, déroule le journal en direct :
chaque ligne écrite par winget, Windows Update ou DISM. Le bouton \menu{Annuler} arrête l'opération
en cours et celles qui attendaient.}

\section{Programmes (winget)}\label{sec:programmes}

\textbf{winget} est le gestionnaire de paquets intégré à Windows 10 et 11 : c'est l'équivalent de
\texttt{yay}. Il connaît plusieurs milliers de logiciels et les applis du Microsoft Store, et sait les
mettre à jour sans fenêtre.

\capture{../capture_programmes-fr.png}

\begin{itemize}
\item Les lignes cochées partent avec \menu{Mettre à jour la sélection}, une par une : l'échec de
l'une n'empêche pas les autres, et le résumé dit lesquelles ont échoué.
\item \menu{Ignorer / ne plus ignorer} retire un logiciel des mises à jour, à la main comme en
automatique (un logiciel dont vous gardez volontairement une ancienne version, par exemple).
\item \menu{Afficher les ignorés} les fait réapparaître, grisés.
\end{itemize}

\attention{Si winget est absent (Windows 10 ancien), installez « Programme d'installation
d'application » depuis le Microsoft Store.}

\section{Hors winget}\label{sec:hors}

Certains logiciels ne sont pas suivis par winget : installés à la main depuis le site de l'éditeur,
ou inconnus de son catalogue. Martinet les répartit en deux listes.

\capture{../capture_hors-fr.png}

\subsection{Logiciels avec une source connue}

Martinet lit la dernière version publiée là où l'éditeur la publie :
\begin{itemize}
\item les \textbf{Releases GitHub} du projet (les préversions « dev », « rc », « beta » sont écartées) ;
\item la \textbf{page de téléchargement} de l'éditeur, lue avec un motif.
\end{itemize}
\menu{Télécharger et installer} récupère l'installeur officiel quand son adresse est connue, puis le
lance : vous suivez l'installation à l'écran. Sinon, \menu{Ouvrir la page} ouvre la page de
téléchargement.

Le catalogue fourni couvre PHD2, ASCOM Platform, N.I.N.A., ASTAP, GraXpert et SharpCap.

\subsection{Logiciels sans source connue}

Pour chacun, trois possibilités :
\begin{enumerate}
\item \menu{Confier à winget…} cherche le logiciel dans le catalogue winget ; s'il y est, winget le
réinstalle par-dessus (les réglages sont en général conservés) et le suit désormais.
\item \menu{Ajouter une source…} indique à Martinet où trouver sa dernière version :
\capture[0.62]{../capture_source-fr.png}
Pour un projet GitHub, il suffit de « propriétaire/dépôt » (l'adresse complète est acceptée). Pour une
page web, le motif de version est une expression régulière dont le premier groupe entre parenthèses
est la version ; la plus haute trouvée sur la page est retenue. \menu{Tester} affiche ce que Martinet
trouve, avant d'enregistrer.
\item \menu{Ignorer} le retire de la liste (pilotes, outils sans mises à jour…).
\end{enumerate}

\astuce{\menu{Sources…} liste le catalogue et vos sources. Modifier une source du catalogue crée
votre propre version, qui la remplace. \menu{Importer…} charge un fichier de sources tout prêt.
Les logiciels qui se mettent à jour eux-mêmes (navigateurs, Microsoft 365, Adobe, NVIDIA App…) sont
marqués « Se met à jour tout seul ».}

\section{Windows et pilotes}\label{sec:windows}

\capture{../capture_windows-fr.png}

\subsection{Carte graphique}

Pour une carte \textbf{NVIDIA}, Martinet interroge le service officiel de NVIDIA (celui de
l'application NVIDIA) avec l'identifiant de la carte, et compare au pilote installé.
\menu{Télécharger et installer le pilote} récupère l'installeur officiel et le lance.

Pour une \textbf{Radeon} (ou les graphismes intégrés d'un Ryzen), Martinet lit la page de pilotes
d'AMD propre à la carte, ou au processeur, et n'en retient que le pilote \textbf{recommandé} (WHQL),
jamais l'« optionnel ». Il le compare à la version d'AMD Software installée. L'installeur d'AMD
s'ouvre et se suit à l'écran : le passage automatique ne fait que le signaler. Les puces Vega et
Polaris sont sur la branche de maintenance d'AMD : nouveaux pilotes rares, Windows Update suffit le
plus souvent. Pour une carte Intel, \menu{Ouvrir la page} mène à l'outil de détection d'Intel.

\subsection{Windows Update}

La liste est celle de \textit{Paramètres > Windows Update}, pilotes des fabricants compris (chipset,
Wi-Fi, audio…). Les mises à jour \textbf{facultatives} (pilotes facultatifs, avant-premières de
Windows) sont incluses tant que la case \menu{Inclure les mises à jour facultatives} est cochée ; ce
choix vaut aussi pour « Mettre à jour et nettoyer » et pour le passage automatique.
\menu{Mettre à jour Defender} rafraîchit les signatures de l'antivirus.

\attention{Après l'installation, Windows demande souvent un redémarrage : Martinet le propose (dans une
minute, le temps d'enregistrer votre travail).}

\section{Raccourcis morts}\label{sec:raccourcis}

\capture{../capture_raccourcis-fr.png}

Martinet lit les raccourcis du Bureau, du menu Démarrer et de la barre des tâches et garde ceux dont la
cible a disparu. Par prudence, il \textbf{ne compte pas} :
\begin{itemize}
\item les raccourcis spéciaux (Panneau de configuration, applis du Store, raccourcis d'installation) ;
\item les cibles sur le réseau ou sur un lecteur absent (clé USB, disque débranché).
\end{itemize}
\menu{Mettre de côté la sélection} les déplace dans les données de Martinet, avec leur emplacement
d'origine. La liste du bas permet de les \menu{Restaurer} à tout moment.

\section{Ménage}\label{sec:menage}

\capture{../capture_menage-fr.png}

Seulement ce que Windows sait recréer : fichiers temporaires de plus de 7 jours, rapports d'erreurs,
cache de téléchargement de Windows Update, installeurs déjà utilisés, vidages mémoire après plantage,
journaux de Windows de plus de 7 jours, restes d'installation des pilotes (\texttt{C:\textbackslash NVIDIA},
\texttt{C:\textbackslash AMD}). Ne sont jamais cochés d'office : le cache des miniatures, les caches des
navigateurs (vidés seulement navigateurs fermés ; ni historique, ni cookies, ni mots de passe), les caches
de shaders (saccades en jeu le temps qu'ils se refassent), les anciennes versions des composants Windows
(DISM, plusieurs minutes) et la \textbf{corbeille}, qui demande confirmation. Un fichier en cours
d'utilisation est laissé en place.

\section{Planification : le mode automatique}\label{sec:planif}

\capture{../capture_planif-fr.png}

Cochez \menu{Passage automatique}, choisissez le jour et l'heure, puis \menu{Enregistrer}. Martinet
crée une tâche planifiée Windows qui :
\begin{itemize}
\item tourne avec vos droits les plus élevés, seulement quand votre session est ouverte ;
\item rattrape au démarrage suivant un passage manqué (ordinateur éteint) ;
\item fait ce qui est coché dans « Ce qui est fait sans vous » : programmes, Windows Update et pilotes,
pilote NVIDIA (installé sans fenêtre), raccourcis morts mis de côté, ménage ;
\item affiche à la fin une \textbf{notification} qui résume le passage ; un clic l'ouvre dans Martinet.
\end{itemize}

Les logiciels hors winget sont seulement \textbf{signalés} : leurs installeurs ne savent pas tous
s'installer sans vous. \menu{Lancer ce passage maintenant} fait le même travail tout de suite.

\section{Historique}

\capture{../capture_historique-fr.png}

Chaque passage, à la main ou automatique, avec le détail ligne par ligne. Le menu
\menu{Fichier > Ouvrir le dossier des données} mène au dossier où Martinet range ses réglages, ses
sources, l'historique et les raccourcis mis de côté (\fichier{\%LOCALAPPDATA\%\textbackslash Martinet}).

\section{Ligne de commande}

\begin{center}
\begin{tabular}{>{\ttfamily}l l}
\toprule
run.py --verifier & liste ce qui est à mettre à jour, sans rien changer \\
run.py --auto & passage complet, comme la tâche planifiée \\
run.py --lang en & interface en anglais \\
\bottomrule
\end{tabular}
\end{center}

\section{En cas de problème}

\begin{itemize}
\item \menu{Aide > Signaler un problème…} envoie votre description avec la version de Martinet et du
système, rien d'autre.
\item Une étape qui échoue (pas de réseau, service Windows Update arrêté…) n'empêche pas les autres :
Martinet dit laquelle, et pourquoi.
\item Un logiciel qui refuse de se mettre à jour par winget est souvent ouvert : fermez-le et relancez.
\end{itemize}

\end{document}
